% Plano de implantação do LibVox
% Compilar com:  pdflatex plano_de_implantacao.tex
%
% Sem babel: esta instalação não traz o portuguese.ldf, e aplicar padrões de
% hifenização ingleses ao português quebra palavra no lugar errado. Por isso a
% hifenização fica desligada, e o texto justificado se apoia na expansão do
% microtype para não abrir vãos entre palavras.
\documentclass[11pt,a4paper]{article}

\usepackage[T1]{fontenc}
\usepackage[utf8]{inputenc}
\usepackage{textcomp}
\usepackage{tgpagella}
\usepackage[scaled=0.92]{helvet}
\usepackage[none]{hyphenat}

\usepackage[a4paper,top=3cm,bottom=3cm,left=3.2cm,right=3.2cm]{geometry}
% Expansão e protrusão fazem o trabalho que a hifenização faria: sem padrões de
% português instalados, é o que mantém o texto justificado sem rios de espaço.
\usepackage[expansion=alltext,protrusion=true,final]{microtype}
\usepackage[parfill]{parskip}
\usepackage{enumitem}
\usepackage{xcolor}
\usepackage{titlesec}
\usepackage[hidelinks]{hyperref}

% Sem babel, \today sai em inglês e o espaçamento após ponto final segue a
% convenção anglófona de espaço duplo. Os dois são corrigidos aqui.
\frenchspacing
\emergencystretch=2em
\newcommand{\mesbr}{\ifcase\month\or janeiro\or fevereiro\or março\or abril\or
  maio\or junho\or julho\or agosto\or setembro\or outubro\or novembro\or
  dezembro\fi}
\newcommand{\databr}{\number\day\ de \mesbr\ de \number\year}

\definecolor{marcia}{HTML}{CF0D4D}
\definecolor{tinta}{HTML}{1B1A21}
\definecolor{suave}{HTML}{55525F}

\color{tinta}

\titleformat{\section}{\sffamily\large\bfseries\color{tinta}}{}{0pt}{}
\titlespacing*{\section}{0pt}{1.6\baselineskip}{0.5\baselineskip}

\setlist[itemize]{leftmargin=1.4em,itemsep=0.45\baselineskip,topsep=0.4\baselineskip}
\setlist[enumerate]{leftmargin=1.9em,itemsep=0.35\baselineskip,topsep=0.4\baselineskip}

% O parskip está com a opção `parfill`, que fixa \parfillskip=30pt plus 1fil:
% todo parágrafo tem de sobrar 30pt na última linha. Um parágrafo cujo conteúdo é
% uma caixa da largura exata da linha (a régua do título, a caixa do pedido)
% estoura por exatamente esses 30pt. Restaurar o valor padrão do LaTeX resolve, e
% precisa ser feito DENTRO do grupo que contém o \par, senão não se aplica.
\newcommand{\caixacheia}{\parfillskip=0pt plus 1fil\relax}

% Rótulo de etapa. As etapas deixaram de ser \section: elas são passos de uma
% só seção, e quatro títulos de seção para quatro passos davam ao que falta mais
% peso na página do que ao que já funciona.
\newcommand{\etapa}[1]{%
  \par\vspace{0.45\baselineskip}%
  \noindent{\sffamily\bfseries #1}\par\vspace{0.1\baselineskip}}

% Caixa do pedido. Sem tcolorbox, que não está no preâmbulo e não vale uma
% dependência nova: \fcolorbox com uma minipage faz o mesmo. A subtração de
% \fboxsep e \fboxrule é o que impede a caixa de passar da margem.
\newcommand{\pedido}[1]{%
  \par\vspace{0.3\baselineskip}%
  {\caixacheia
  \noindent\fcolorbox{marcia}{marcia!5}{%
    \begin{minipage}{\dimexpr\linewidth-2\fboxsep-2\fboxrule\relax}%
      \vspace{0.2\baselineskip}#1\vspace{0.2\baselineskip}%
    \end{minipage}}\par}%
  \vspace{0.3\baselineskip}}

\newcommand{\titulo}{%
  \noindent{\sffamily\LARGE\bfseries Plano de implantação do LibVox}\par
  \vspace{0.35em}
  {\caixacheia\color{marcia}\rule{\linewidth}{1.2pt}\par}
  \vspace{0.4em}
  {\sffamily\small\color{suave} IMPA Tech \quad\textbullet\quad \databr}\par
  \vspace{1.6\baselineskip}}

\begin{document}
\titulo

O LibVox é o nome de um projeto de Iniciação Científica (IC) do IMPA Tech. Nele, implementamos uma assistente conversacional baseada em linguagem natural para catálogos de bibliotecas, chamada \textbf{MARC.IA}. Com ela, o leitor pode fazer perguntas como: \emph{``vocês têm alguma coisa sobre xadrez para iniciantes?''}, e recebe como resposta os títulos disponíveis no catálogo, com número de exemplares e localização.

\section{O que já implementamos?}

O LibVox já responde hoje, de ponta a ponta: a pergunta escrita em linguagem
natural é traduzida em consultas ao acervo e volta como uma resposta com título,
autor, quantidade de exemplares e localização.

\begin{itemize}
  % PREENCHER: quantas bibliotecas estão configuradas, quantos registros somam os
  % acervos carregados, e desde quando o ambiente de testes está no ar. Estes três
  % números valem mais, nesta seção, do que qualquer adjetivo.
  \item Há um ambiente de testes no ar, com acervos reais e uma página
        de demonstração.
  \item Uma única instalação pode atender várias bibliotecas ao mesmo tempo.
  \item O código é aberto, sob licença AGPL-3.0, com o copyright do IMPA Tech.
  \item Alunos do IMPA Tech já trabalham no projeto, pegando tarefas do
        desenvolvimento e planejando a escrita de artigos científicos. 
\end{itemize}

\section{Por que isso interessa ao IMPA Tech}

Nosso objetivo é manter o LibVox como projeto permanente do IMPA Tech.

\begin{enumerate}
  \item \textbf{Ensino e extensão no mesmo trabalho.} Alunos escrevem código
        que vai para produção em uma biblioteca real.
  \item \textbf{Alcance} Sendo código aberto, qualquer
        instituição pode subir a própria instância, o que leva o nome do 
        IMPA Tech para bibliotecas de todo o Brasil.
  \item \textbf{Produção científica.} Há uma agenda de pesquisa com diversas linhas mapeadas.
\end{enumerate}

\section{O que falta para entrar em produção}

\etapa{Etapa 1. Ambiente de testes}
\begin{itemize}
  \item Configurar uma máquina dentro do datacenter do IMPA para prover o modelo
        de linguagem que interpreta as perguntas e indica quais consultas o
        LibVox deve fazer (Máquina com GPU). Esse recurso será utilizado apenas para atender as demandas da nossa biblioteca.
        {\color{suave}(Analisando com o Daniel e Lucas)}
  \item Criar uma instância do Koha da nossa biblioteca dedicada a testes, com
        cópia do acervo.
        {\color{suave}(Analisando com Sérgio Piloto)}
  % A VM é externa porque o datacenter do IMPA não tem como oferecer esse recurso
  % agora. Redigido sem apontar o datacenter como obstáculo: a frase afirma o que
  % resolve e deixa a porta aberta para internalizar depois.
  \item Contratar uma máquina virtual para hospedar os componentes do LibVox: o
        núcleo da assistente, o portal de administração e o banco de dados. O
        datacenter do IMPA pode assumir esse papel mais adiante; para começar,
        uma máquina contratada resolve.
        {\color{suave}(Custo Anual \$100 dólares)}
\end{itemize}

% Etapas 2, 3 e 4 têm uma ação cada. Lista de um item só gasta espaço vertical
% e não organiza nada: vão como parágrafo.
\etapa{Etapa 2. Acesso ao catálogo}
Criar um usuário de API no Koha com \textbf{permissão apenas de consulta ao
catálogo}.

\etapa{Etapa 3. Alteração no catálogo}
Acrescentar uma linha ao template da tela de consulta ao acervo, que faz aparecer
um botão no canto da página.

\etapa{Etapa 4. Política de uso}
Redigir o texto apresentado ao leitor ao clicar no botão, antes da primeira
pergunta. Ele precisa informar o que é a ferramenta, que o acesso é anônimo, que
as perguntas e respostas são guardadas, com quais finalidades, e que o leitor não
deve escrever dados pessoais na conversa.

\section{Mapeamento de Riscos}

\begin{itemize}
  \item \textbf{O acesso ao catálogo é somente de leitura.} É o próprio Koha que
        recusa qualquer outra operação. Não há acesso a dados de usuários,
        empréstimos, reservas ou multas, nem qualquer função de edição.
  \item \textbf{O sistema guarda apenas a pergunta e a resposta.} O desenho do
        banco de dados é o que garante isso: ele simplesmente não tem campo para
        nome, e-mail, telefone, etc.
  \item \textbf{As buscas não chegam ao servidor da biblioteca.} O LibVox
        trabalha sobre uma cópia do acervo, que mantém atualizada. A única
        consulta em tempo real é a de disponibilidade de exemplar.
  \item \textbf{A alteração no catálogo é uma linha, e é reversível.} É a única
        modificação no Koha: nenhum registro, configuração ou comportamento da
        busca é alterado.
\end{itemize}

\section{Como funciona o LibVox e tratamento de alucinações}

Duas características do desenho mantêm a resposta ancorada no acervo, e são elas
que reduzem o risco de alucinação. O modelo pode variar de uma instância para
outra: na nossa, usaremos o Qwen, que é um modelo aberto.

\begin{enumerate}
  \item \textbf{Todo título vem do catálogo.} O modelo de linguagem escreve a
        resposta para o usuário utilizando as informações do acervo da
        biblioteca. Portanto, ele não inventa títulos de livros ou informações
        que não estejam no catálogo.

  \item \textbf{A resposta diz o que foi buscado.} A assistente informa qual
        termo pesquisou no catálogo. Com isso o leitor pode perceber quando a
        busca não correspondeu ao que ele pediu, e refazer a pergunta com outras
        palavras.
\end{enumerate}

Além da prevenção, há acompanhamento: as conversas ficam gravadas e serão lidas
periodicamente para verificar a qualidade das respostas. Servem a dois fins,
declarados desde já: melhorar a ferramenta e sustentar publicações futuras.
\textbf{O registro é integralmente anônimo}, pela mesma razão da seção anterior:
não existe no sistema onde guardar dado que identifique quem perguntou.

\section{Manutenção depois de aberta ao público}

Esta primeira instância atende apenas os usuários da biblioteca do IMPA, e por
isso a manutenção é feita pela própria equipe. São duas tarefas: reiniciar algum
serviço caso ele trave, e atualizar o sistema quando os alunos produzirem
novidade.


\section{O que pedimos}

\pedido{%
\begin{enumerate}[leftmargin=1.6em,topsep=0pt,itemsep=0.3\baselineskip]
  \item \textbf{A máquina para o modelo de linguagem, no datacenter, e a
        instância do Koha de testes.} São os dois recursos internos que a Etapa 1
        depende.
  \item \textbf{Uma máquina virtual externa para hospedar os componentes do
        LibVox.} É onde a assistente, o portal e o banco de dados rodam. O
        datacenter pode assumir esse papel mais adiante; para entrar em
        produção agora, uma máquina contratada resolve.
\end{enumerate}}

\end{document}
